\chapter{Vorher, nachher, und mit wem}

\section{Die doppelte Differenz}

\begin{merksatz}
\boxtitel{Eine Reform in manchen Kassen}
Zum Jahr $3$ führen einige Kassen eine Zuzahlung ein, andere nicht. Man beobachtet beide Gruppen vorher und nachher.

\begin{center}
{\small
\renewcommand{\arraystretch}{1.5}
\begin{tabular}{@{}P{2.4cm}P{2.0cm}P{2.0cm}P{2.2cm}@{}}
\toprule
& \textbf{vorher} & \textbf{nachher} & \textbf{Differenz}\\
\midrule
Reformkassen & $\bar y_{R,0}$ & $\bar y_{R,1}$ & $\Delta_R$\\
andere Kassen & $\bar y_{K,0}$ & $\bar y_{K,1}$ & $\Delta_K$\\
\midrule
& & & $\Delta_R - \Delta_K$\\
\bottomrule
\end{tabular}}
\end{center}

\begin{center}
\begin{minipage}{0.9\textwidth}
\textbf{Die Veränderung der Betroffenen, abzüglich der Veränderung der Unbetroffenen.}
\end{minipage}
\end{center}

\begin{center}
\begin{minipage}{0.88\textwidth}
\kursiv{\textbf{Band 23, Kapitel~10 hat die doppelte Differenz als natürlichen Versuch eingeführt}; hier wird sie zum Panelmodell. Die Rechnung ist dieselbe: Die erste Differenz beseitigt alles Zeitkonstante der Gruppen, die zweite alles, was beide Gruppen gleichermaßen trifft. Was bleibt, ist die Reform --- wenn eine Annahme gilt.}
\end{minipage}
\end{center}
\end{merksatz}

\begin{definition}
\boxtitelm{Definition 9.1 --- parallele Trends}
Ohne Reform hätten sich beide Gruppen gleich entwickelt:
\[ \mathbb{E}\bigl[y_{1}(0) - y_{0}(0)\,\big|\,R\bigr] = \mathbb{E}\bigl[y_{1}(0) - y_{0}(0)\,\big|\,K\bigr] . \]
Dabei ist $y_t(0)$ das Ergebnis ohne Reform --- die potentielle Größe aus Band 23.
\end{definition}

\begin{merksatz}
\boxtitel{Satz 9.2 --- was die doppelte Differenz schätzt}
Unter parallelen Trends gilt
\[ \Delta_R - \Delta_K \ \longrightarrow\ \mathbb{E}\bigl[y_1(1) - y_1(0)\,\big|\,R\bigr] \]
--- der mittlere Effekt der Reform bei den Betroffenen.

\schritt{Beweis}
\[ \Delta_R = \mathbb{E}[y_1(1)-y_0(0)\mid R] = \underbrace{\mathbb{E}[y_1(1)-y_1(0)\mid R]}_{\text{Effekt}} + \underbrace{\mathbb{E}[y_1(0)-y_0(0)\mid R]}_{\text{Trend der Betroffenen}} \]
\[ \Delta_K = \mathbb{E}[y_1(0)-y_0(0)\mid K] \]
Nach Definition 9.1 sind die beiden Trends gleich; die Differenz lässt den Effekt übrig. $\blacksquare$

\begin{center}
\begin{minipage}{0.9\textwidth}
\textbf{Die Gruppen dürfen beliebig verschieden sein. Sie müssen sich nur gleich entwickeln.}
\end{minipage}
\end{center}

\schritt{Als Regression}
\[ y_{it} = \alpha_i + \delta_t + \tau\,D_{it} + u_{it}, \qquad D_{it} = \mathbf 1\{i \in R,\ t \ge 3\} \]
\kursiv{Das zweiseitige Modell aus Kapitel~7 mit einem Behandlungsindikator --- bei zwei Gruppen und zwei Perioden ist $\hat\tau$ genau $\Delta_R-\Delta_K$.}
\end{merksatz}

\section{Die Annahme prüfen}

\begin{merksatz}
\boxtitel{Vorperioden}
Parallele Trends sind nicht prüfbar --- sie betreffen eine Welt ohne Reform. Prüfbar ist, ob die Gruppen sich \kursiv{vor} der Reform parallel entwickelt haben.

\begin{center}
{\small
\renewcommand{\arraystretch}{1.48}
\begin{tabular}{@{}P{3.6cm}P{2.4cm}P{2.0cm}@{}}
\toprule
\textbf{Vergleich} & \textbf{Schätzung} & \textbf{Soll}\\
\midrule
Jahr $2$ gegen $1$ (vorher) & $-0{,}04$ & $0$\\
Jahr $3$ gegen $1$--$2$ & $0{,}92$ & $1$\\
Jahr $5$ gegen $1$--$2$ & $3{,}03$ & $3$\\
\bottomrule
\end{tabular}}
\end{center}
\kursiv{Simuliert: Reform in Jahr $3$, Wirkung wächst um eins je Jahr.}

\begin{center}
\begin{minipage}{0.88\textwidth}
\kursiv{\textbf{Die erste Zeile ist der Vorperiodentest, und sein Bestehen ist beruhigend, aber kein Beweis}: Dass die Trends vorher parallel waren, sagt nichts darüber, ob sie es ohne Reform geblieben wären --- etwa wenn die Reform gerade deshalb kam, weil sich etwas zu ändern begann. Band 23 nannte das die Ashenfelter-Delle. Die zweite und dritte Zeile zeigen die Ereignisstudie: der Effekt Jahr für Jahr, statt als eine Zahl.}
\end{minipage}
\end{center}
\end{merksatz}

\section{Der Bruchfall}

\begin{warnung}
\boxtitelr{Gestaffelte Einführung}
Nun führe eine dritte Gruppe die Zuzahlung erst im Jahr $6$ ein. Man schätzt das zweiseitige Modell über alle:

\begin{center}
{\small
\renewcommand{\arraystretch}{1.48}
\begin{tabular}{@{}P{4.2cm}P{2.4cm}@{}}
\toprule
\textbf{Schätzer} & \textbf{Effekt}\\
\midrule
zweiseitige feste Effekte & $1{,}82$\\
Mittel sauberer Vergleiche & $2{,}98$\\
\midrule
wahrer mittlerer Effekt & $\mathbf{2{,}99}$\\
\bottomrule
\end{tabular}}
\end{center}

\begin{center}
\begin{minipage}{0.9\textwidth}
\textbf{Das Standardverfahren unterschätzt um vierzig Prozent --- bei korrekt spezifizierten festen Effekten und erfüllten parallelen Trends.}
\end{minipage}
\end{center}

\schritt{Woran es liegt}
\begin{center}
\begin{minipage}{0.88\textwidth}
\kursiv{Das zweiseitige Modell benutzt \kursiv{alle} Vergleiche zwischen Gruppen und Zeiten --- auch den, bei dem die Spätgruppe behandelt wird und die Frühgruppe als Kontrolle dient. \textbf{Die Frühgruppe ist aber zu diesem Zeitpunkt selbst behandelt, und ihr Effekt wächst noch.} Dieser \enquote{verbotene Vergleich} zieht den wachsenden Effekt der Frühgruppe vom Effekt der Spätgruppe ab:}
\end{minipage}
\end{center}

\begin{center}
{\small
\renewcommand{\arraystretch}{1.48}
\begin{tabular}{@{}P{4.4cm}P{2.0cm}P{1.6cm}@{}}
\toprule
\textbf{Vergleich} & \textbf{Schätzung} & \textbf{wahr}\\
\midrule
spät gegen nie, Jahr $7$ & $2{,}0$ & $2$\\
spät gegen \kursiv{früh}, Jahr $7$ gegen $5$ & $\mathbf{0{,}00}$ & $2$\\
\bottomrule
\end{tabular}}
\end{center}

\begin{center}
\begin{minipage}{0.88\textwidth}
\kursiv{\textbf{Der verbotene Vergleich liefert null}, weil die Frühgruppe zwischen Jahr $5$ und $7$ selbst von Effekt $3$ auf $5$ wächst --- genau so viel wie die Spätgruppe gewinnt. Das zweiseitige Modell mischt solche Vergleiche mit negativem Gewicht ein. \textbf{Bei konstantem Effekt stört das nicht} (die Simulation ergibt dann $2{,}04$ bei wahrem Wert $2$); sobald die Wirkung über die Zeit wächst oder zwischen Gruppen differiert, schon.}
\end{minipage}
\end{center}

\schritt{Die Abhilfe}
\begin{center}
\begin{minipage}{0.88\textwidth}
\kursiv{Man rechnet nur saubere Zweiergruppenvergleiche --- jede Behandlungsgruppe gegen die Nie-Behandelten oder gegen die Noch-nicht-Behandelten --- und mittelt sie mit sinnvollen Gewichten. Das ergibt $2{,}98$. \textbf{Die Verfahren dafür sind seit 2020 Standard}, und sie haben gezeigt, dass viele ältere Befunde zu klein, manche mit falschem Vorzeichen geschätzt waren.}
\end{minipage}
\end{center}
\end{warnung}

\section{Übungsaufgaben}

\begin{aufgabe}
\boxtitel{Aufgabe 9.1 --- Die Tabelle ausfüllen}
Reformkassen: $7{,}2$ Besuche vorher, $6{,}1$ nachher. Andere Kassen: $6{,}8$ vorher, $6{,}5$ nachher. Die Reform erhöhte die Zuzahlung um $8$ Euro.

\loesung
\schritt{Schritt 1 --- Die beiden Differenzen}
\[ \Delta_R = 6{,}1-7{,}2 = -1{,}1, \qquad \Delta_K = 6{,}5-6{,}8 = -0{,}3 \]

\schritt{Schritt 2 --- Die doppelte}
\[ \hat\tau = -1{,}1-(-0{,}3) = \mathbf{-0{,}8} \ \text{Besuche} \]

\schritt{Schritt 3 --- Je Euro}
\[ \frac{-0{,}8}{8} = -0{,}100 \]

\schritt{Schritt 4 --- Was die einfachen Vergleiche gesagt hätten}
\begin{center}
{\small
\renewcommand{\arraystretch}{1.48}
\begin{tabular}{@{}P{3.6cm}P{2.0cm}P{2.6cm}@{}}
\toprule
\textbf{Vergleich} & \textbf{Effekt} & \textbf{Fehler}\\
\midrule
nur vorher/nachher & $-1{,}1$ & enthält den Trend\\
nur Reform/andere, nachher & $-0{,}4$ & enthält das Niveau\\
doppelt & $\mathbf{-0{,}8}$ & ---\\
\bottomrule
\end{tabular}}
\end{center}

\begin{center}
\begin{minipage}{0.88\textwidth}
\kursiv{\textbf{Der reine Vorher-nachher-Vergleich überschätzt, weil die Besuche ohnehin zurückgingen; der reine Gruppenvergleich unterschätzt, weil die Reformkassen schon vorher höhere Werte hatten.} Beide Fehler sind typisch, und beide Verfahren sind in Evaluationsberichten verbreitet. Die doppelte Differenz braucht nur vier Zahlen mehr.}
\end{minipage}
\end{center}
\end{aufgabe}

\begin{aufgabe}
\boxtitel{Aufgabe 9.2 --- Wenn die Trends nicht parallel sind}
Die Reformkassen hatten schon vor der Reform einen stärkeren Rückgang: $0{,}2$ Besuche je Jahr mehr als die anderen. Was folgt?

\loesung
\schritt{Schritt 1 --- Die Verzerrung}
Satz 9.2 liefert dann Effekt plus Trenddifferenz:
\[ \hat\tau = \tau + (-0{,}2) \quad\Longrightarrow\quad \tau = -0{,}8+0{,}2 = -0{,}6 \]

\schritt{Schritt 2 --- Die Abhilfe}
\begin{center}
{\small
\renewcommand{\arraystretch}{1.5}
\begin{tabular}{@{}P{3.0cm}P{5.6cm}@{}}
\toprule
\textbf{Trendbereinigung} & Man schätzt den Vortrend aus den Vorperioden und zieht ihn ab --- verlangt mehrere Vorperioden.\\
\textbf{Dreifache Differenz} & Eine weitere Gruppe, die den Trend teilt, aber nicht die Reform --- etwa Versicherte derselben Kassen, die von der Zuzahlung befreit sind.\\
\textbf{Schranken} & Man gibt an, wie groß die Trendabweichung höchstens sein darf, damit das Vorzeichen erhalten bleibt.\\
\bottomrule
\end{tabular}}
\end{center}

\schritt{Schritt 3 --- Die Schranke im Beispiel}
\begin{center}
\begin{minipage}{0.9\textwidth}
\textbf{Der Effekt wird erst bei einer Trendabweichung von $0{,}8$ Besuchen je Jahr zu null --- dem Vierfachen der beobachteten. Das Ergebnis ist robust.}
\end{minipage}
\end{center}

\begin{center}
\begin{minipage}{0.88\textwidth}
\kursiv{\textbf{Das ist die Empfindlichkeitsanalyse aus Band 23, Kapitel~7, auf die Trendannahme angewandt.} Statt zu behaupten, die Trends seien parallel, sagt man, wie unparallel sie sein dürften --- und ob eine solche Abweichung plausibel ist. Seit 2023 gibt es dafür ausgearbeitete Verfahren, und sie gehören inzwischen zum Standard guter Ereignisstudien.}
\end{minipage}
\end{center}
\end{aufgabe}

\begin{aufgabe}
\boxtitel{Aufgabe 9.3 --- Die Gewichte des zweiseitigen Modells}
Zeigen Sie an einem Minimalbeispiel, dass das zweiseitige Modell einen verbotenen Vergleich enthält.

\loesung
\schritt{Schritt 1 --- Zwei Gruppen, drei Perioden}
Gruppe~A behandelt ab Periode $2$, Gruppe~B ab Periode $3$, keine Kontrollgruppe.

\schritt{Schritt 2 --- Die möglichen Zweiervergleiche}
\begin{center}
{\small
\renewcommand{\arraystretch}{1.5}
\begin{tabular}{@{}P{3.2cm}P{2.4cm}P{2.6cm}@{}}
\toprule
\textbf{Vergleich} & \textbf{Perioden} & \textbf{Kontrolle}\\
\midrule
A gegen B & $1 \to 2$ & B noch unbehandelt \checkmark\\
B gegen A & $2 \to 3$ & A \kursiv{schon behandelt} \textcolor{akzent}{$\times$}\\
\bottomrule
\end{tabular}}
\end{center}

\schritt{Schritt 3 --- Was im zweiten Vergleich passiert}
Der Effekt von A sei $\tau_{A,1}$ in der ersten und $\tau_{A,2}$ in der zweiten Behandlungsperiode:
\[ \text{Vergleich 2} = \tau_{B,1} - (\tau_{A,2}-\tau_{A,1}) \]

\begin{center}
\begin{minipage}{0.9\textwidth}
\textbf{Wächst der Effekt von A, wird er vom Effekt von B abgezogen. Das zweiseitige Modell mittelt beide Vergleiche --- und der zweite ist falsch.}
\end{minipage}
\end{center}

\schritt{Schritt 4 --- Die Bedingung für Unschädlichkeit}
\begin{center}
\begin{minipage}{0.88\textwidth}
\kursiv{Nur wenn $\tau_{A,2} = \tau_{A,1}$ --- der Effekt also über die Zeit konstant ist --- verschwindet die Störung. \textbf{Genau das ist bei Zuzahlungen unwahrscheinlich}: Patienten lernen, Ärzte passen sich an, Termine verschieben sich. Wer gestaffelte Einführungen mit dem zweiseitigen Modell auswertet, nimmt stillschweigend an, dass nichts davon geschieht.}
\end{minipage}
\end{center}
\end{aufgabe}

\begin{geschichte}
\boxtitelg{Cholera, Mindestlöhne und ein Aufruhr}
\Person{John Snow} benutzte 1855 eine doppelte Differenz, lange bevor es den Namen gab: Londoner Bezirke, die ihr Wasser von zwei Gesellschaften bezogen, verglichen vor und nach dem Umzug einer Entnahmestelle flussaufwärts --- Band 23 erzählt es.

\begin{center}
\begin{minipage}{0.9\textwidth}
\kursiv{\Person{David Card} und \Person{Alan Krueger} machten das Verfahren 1994 berühmt: Sie verglichen Schnellrestaurants in New Jersey, wo der Mindestlohn stieg, mit solchen in Pennsylvania, wo er gleich blieb --- und fanden keinen Beschäftigungsrückgang. \textbf{Der Befund widersprach dem Lehrbuch und löste einen Streit aus, der bis heute anhält}; Card erhielt 2021 den Nobelpreis.}
\end{minipage}
\end{center}

\schritt{Die Krise des Verfahrens}
\begin{center}
\begin{minipage}{0.9\textwidth}
\kursiv{Zwischen 2018 und 2021 zeigten \Person{Goodman-Bacon}, \Person{de Chaisemartin} und \Person{D'Haultf\oe uille}, \Person{Callaway} und \Person{Sant'Anna} sowie \Person{Sun} und \Person{Abraham} unabhängig voneinander den Bruchfall dieses Kapitels. \textbf{Innerhalb weniger Jahre musste die angewandte Ökonomie eines ihrer meistbenutzten Werkzeuge neu lernen} --- ein seltener Fall, in dem eine methodische Einsicht eine ganze Literatur zur Nachprüfung zwang.}
\end{minipage}
\end{center}

\schritt{Und Deutschland}
\begin{center}
\begin{minipage}{0.9\textwidth}
\kursiv{Die Praxisgebühr von 2004 bis 2012 --- zehn Euro je Quartal --- ist der deutsche Leitfall für dieses Kapitel. \textbf{Studien mit dem Sozio-oekonomischen Panel fanden einen Rückgang der Arztbesuche um etwa zehn Prozent, konzentriert auf Geringverdiener}; andere fanden kaum etwas. Die Unterschiede gehen auf die Wahl der Kontrollgruppe zurück, denn befreit waren nur Kinder --- und Kinder sind als Vergleich für Erwachsene fragwürdig.}
\end{minipage}
\end{center}
\end{geschichte}

\begin{forschung}
\boxtitelp{Doppelte Differenz heute}
\begin{center}
{\small
\renewcommand{\arraystretch}{1.45}
\begin{tabular}{@{}P{3.2cm}P{5.0cm}@{}}
\toprule
\textbf{Gestaffelte Verfahren} & Robuste Schätzer für heterogene, wachsende Effekte.\\
\textbf{Empfindlichkeit} & Schranken für Abweichungen von parallelen Trends.\\
\textbf{Synthetische Kontrolle} & Eine Kontrollgruppe aus gewichteten Nichtbetroffenen bauen.\\
\bottomrule
\end{tabular}}
\end{center}

\schritt{Zur dritten Zeile}
\begin{center}
\begin{minipage}{0.88\textwidth}
\kursiv{Gibt es nur eine behandelte Einheit --- ein Bundesland, eine Klinik ---, so wählt man Gewichte für die anderen so, dass die gewichtete Summe den Vorverlauf der behandelten nachbildet, und schreibt die Abweichung danach der Maßnahme zu. \textbf{Das Verfahren von \Person{Abadie} aus dem Jahr 2003 ist inzwischen das wichtigste Werkzeug für Fallstudien}, und seine Verbindung zur doppelten Differenz ist seit 2021 gut verstanden: Es ist eine doppelte Differenz mit datengetriebener Kontrollgruppe.}
\end{minipage}
\end{center}
\end{forschung}

\begin{merksatz}
\boxtitel{Was Sie jetzt können}
\begin{center}
\renewcommand{\arraystretch}{1.5}
\begin{tabular}{@{}P{3.0cm}P{5.6cm}@{}}
\toprule
\textbf{Definition 9.1} & parallele Trends, nicht gleiche Niveaus\\
\textbf{Satz 9.2} & Effekt bei den Betroffenen\\
\textbf{Vorperioden} & prüfbar, aber kein Beweis\\
\textbf{Der Bruchfall} & verbotene Vergleiche bei Staffelung\\
\textbf{Schranken} & wie unparallel darf es sein\\
\bottomrule
\end{tabular}
\end{center}

\begin{center}
\begin{minipage}{0.88\textwidth}
\kursiv{Damit endet Teil~III. Alle bisherigen Daten hatten viele Einheiten und wenige Zeitpunkte. Teil~IV kehrt das um: eine Reihe, viele Zeitpunkte --- und ein Problem, das die Querschnittsökonometrie nicht kennt.}
\end{minipage}
\end{center}
\end{merksatz}
